\section{The framework}\label{the-framework}

A world is a table of stages; world 1's has 95 rows. A \textbf{node} is one stage for one place (\texttt{assembly@hangar}) or for the world, and one runner makes, judges and accepts every node. The creator's picks (concepts, style, references) are the inputs; from there everything runs without a typed yes, and the creator reviews.

\subsection{Keys, gates and accepts}\label{keys-gates-and-accepts}

\textbf{Keys.} A node is made only when its key has no finish record. The key hashes the stage and place, a code version, the contracts of its output and inputs, its parameters, a seed derived from the world's, the pinned model images, every world file it reads and the accepted output of every node it reads. The code version hashes every source file the stage imports, leaving out the runner and the renting code, which change what a run costs, not what it makes. A fault fixed in one stage moves only its key and those after it. Outputs are kept by content hash, never overwritten.

\textbf{Gates and accepts.} Each gate answers pass, fail or unknown, with measures and a reason; missing evidence is unknown, and unknown blocks like fail. The runner accepts an output by itself when every gate passed and it is no worse than the one accepted before: no answer dropped, no measure moved the wrong way, no gate missing. Later nodes read only accepted outputs. The creator keeps four commands, from his own terminal: reject an accepted output (it moves back to the one before, and every node that read it is stale), accept a worse one with his reason, waive one gate for one key, and override the runner for one paid run. None is needed for a run; a barrier keeps agents from them (a terminal, a store that refuses the creator's records from an agent's process, and a hook).

\textbf{Spending and review.} A paid node is cleared only when every gate on its inputs passed and the month's ceiling has room for its estimate (Section 3.7); everything that rents a machine or calls a paid model refuses without that ticket, and buying again more than EUR 5 of accepted outputs needs a written reason. Each node's review page is built from the store alone, and the creator reviews a place in a USD viewer on the store's copy of its accepted assembly; nothing is rendered for review.

\subsection{From concept to scene}\label{from-concept-to-scene}

\textbf{Inventory, plan and record.} A coding agent lists every visible element of the concept as an inventory row (kind, anchor, size, count) or drops it with a written reason; nothing is made or placed that is not a row. It draws a dimensioned plan, gated on \texttt{mapped} (every element a row, a door or a written drop; every door says where it leads). The scene record is written from the plan, inventory, made ground and sky, gated on \texttt{game-free} (nothing copied from the game the world came from), \texttt{plan-matched} (the shell within 5 cm of the plan), \texttt{openings} (every doorway onto a floor, no view into nothing) and \texttt{seed-matched} (every old entry matched or dropped with a reason).

\textbf{The canonical scene.} Every place is one OpenUSD stage \citep{openusd}; each object, one per real-world object, carries its kind, transform, collision and each part's library surface. Its layers lie in a fixed order, the creator's edit layer on top, then people, sound, light, motion, effects, crops, machines, settle, ground and the generated base, so an edit survives a regenerated base. Each place is assembled once, from every layer's accepted output; its \texttt{layers} gate answers unknown while a needed layer has none, and an old output never stands in.

\textbf{Models, generation first.} Every object goes the prop route first: Nano Banana Pro \citep{nanobananapro} draws its close-up from its crop of the concept, Pixal3D \citep{li2026pixal3d} builds it, the surface stage paints it, and a likeness gate compares its outline from eight sides with its close-up. Only an object whose checks fail is built by a code builder a coding agent writes, reversing the earlier parts check (the creator's choice, 12 October 2026). On 13 of the hub's objects both orders end the same (Table 1), but generation first draws a close-up of every object. Every close-up of eight thin and flat objects failed the shape check, so code builds carry them. The hatch wheel shows the gain: its code build's outline overlap was 0.45, a fail, and its generated model's 0.60. Most of world 1's models predate the change and enter as unverified seeds.

\par\penalty-100\ifdim\dimexpr\pagegoal-\pagetotal\relax<12\baselineskip\clearpage\fi

\textbf{Table 1.} How 13 hub objects, chosen to cover the kinds, end in each order. Close-ups cost EUR 8.12 and the sample's prop route EUR 13.80, about EUR 2.20 an object; no take was baked (the run stopped at the keyboard's paint gate).

\begingroup\footnotesize\setlength{\tabcolsep}{3pt}

{\def\LTcaptype{none} % do not increment counter
\begin{longtable}[]{@{}llll@{}}
\toprule\noalign{}
Order & Prop route & Code build & Needs a new builder \\
\midrule\noalign{}
\endhead
\bottomrule\noalign{}
\endlastfoot
Generation first & 4 & 7 & 2 \\
Code first & 4 & 7 & 2 \\
\end{longtable}
}

\endgroup

\begin{center}\fbox{\begin{minipage}{0.94\linewidth}\setlength{\parskip}{4pt}\small

\textbf{Gap: Figures 2 to 5.} (2) The stages on the hub, from concept to assembled room; (3) the greenhouse's review page; (4) the wreck with an edit that survives a regenerated base whose models were repainted; (5) four of Table 1's objects from close-up to final model. They wait on those places' assemblies and on renders of the sample's takes; the earlier hub figure predates the runner and is withdrawn, and the comparison recorded no agent hours.

\end{minipage}}\end{center}
